\section{Experiments Continued}
    \label{sec:exp_con}

\noindent 
\textbf{Segmentation evaluation.} For a fair comparison, the segmentation metrics (mDice, mHD95) reported in~\cref{tab:lv_segmentation_comparison} are evaluated exclusively on the annotated ED and ES frames, consistent with the sparse supervision protocol. 

\noindent
\textbf{Data and calibration.} All results for the medical metrics in~\cref{tab:lv_segmentation_comparison} are derived from the official CAMUS test subset.
Notably, we report the raw correlation, bias, and standard deviation without applying any bias-correction or post-calibration techniques to reflect the model's direct predictive performance.

\noindent
\textbf{Training dynamics and convergence.}
As illustrated in~\cref{fig:training_progression}, \mymodel exhibits a clear transition in both segmentation accuracy and feature localization.
In the initial phase (iterations 10--30), the segmentation masks are characterized by coarse boundaries, while the corresponding feature activation maps~(\cref{fig:training_iter10,fig:training_iter20,fig:training_iter30}, bottom rows) show diffused activations across the spatial domain.
By iteration 1250 (epoch 19), the model produces spatially precise segmentations~(\cref{fig:training_iter1250}), with feature activations concentrated strictly on the target anatomical structures.
This progression indicates that the training protocol facilitates the transition from global spatial search to fine-grained structural delineation without performance degradation.

\noindent
\textbf{Temporal consistency analysis.}
The visualization across the frame sequences in~\cref{fig:training_progression} demonstrates the capacity of the model to maintain temporal stability.
In the early stages, predictions are temporally inconsistent, with visible fluctuations in mask geometry across frames.
However, the converged model~(\cref{fig:training_iter1250}) preserves structural identity and boundary continuity throughout the cardiac cycle.
This stability is attributed to the following mechanisms.

\noindent
\textbf{Spatiotemporal feature aggregation.} The integration of contextual information from the temporal sequence facilitates the reduction of frame-wise prediction noise by utilizing the temporal evolution of anatomical structures.

\noindent
\textbf{Gradient propagation.} Backpropagating the loss through the temporal network enforces a consistency constraint on intermediate frames, ensuring that non-annotated frames align with the patterns learned from supervised frames.

\noindent
\textbf{Sequential context integration.} The refinement of internal representations as the sequence progresses allows the model to utilize cumulative temporal evidence for more accurate boundary localization, effectively correcting initial segmentation noise.

As shown in~\cref{fig:training_progression}, the segmentation quality of intermediate frames without direct supervision remains consistent with the supervised frames.
This indicates that the architecture effectively enforces temporal consistency and preserves structural identity throughout the sequence.